\begin{remark}\label{sga3:XXI:6.7.8}
It is immediate that \(\operatorname{Aut}(\cR)\) (respectively
\(E_\Delta(\cR)\)) is finite if and only if
\[
 \rgred(\cR)-\rgss(\cR)\leq1.
\]
\end{remark}

\subsection{\(p\)-morphisms of reduced root data}
\label{sga3:XXI:6.8}

Throughout this subsection, \(p\) is a fixed positive integer.

\begin{definition}\label{sga3:XXI:6.8.1}
Let
\[
 \cR=(M,M^\ast,R,R^\ast),
 \qquad
 \cR'=(M',M'^\ast,R',R'^\ast)
\]
be two reduced root data. A group homomorphism
\[
 f:M'\longrightarrow M
\]
is called a \emph{\(p\)-morphism from \(\cR'\) to \(\cR\)} if
there exist a bijection
\[
 u:R\xrightarrow{\sim}R'
\]
and a map
\[
 q:R\longrightarrow\{p^n\mid n\in\mathbb N\}
\]
such that:
\begin{enumerate}
\item[(i)] \(f(u(\alpha))=q(\alpha)\alpha\) for every
  \(\alpha\in R\);
\item[(ii)] \({}^tf(\alpha^\ast)=q(\alpha)u(\alpha)^\ast\) for
  every \(\alpha\in R\).\footnote{Editors' note: These two
  conditions imply
  \[
   q(\alpha)\bigl(u(\alpha)^\ast,u(\beta)\bigr)
   =q(\beta)(\alpha^\ast,\beta)
  \]
  for all \(\alpha,\beta\in R\).}
\end{enumerate}
\end{definition}

\begin{corollary}\label{sga3:XXI:6.8.2}
A \(1\)-morphism is precisely a morphism.
\end{corollary}

\begin{corollary}\label{sga3:XXI:6.8.3}
The transpose of a \(p\)-morphism is a \(p\)-morphism.
\end{corollary}

\begin{lemma}\label{sga3:XXI:6.8.4}
If \(w\in W(\cR)\) and \(\alpha\in R\), then
\[
 q(w(\alpha))=q(\alpha).
\]
The map
\[
 s_\alpha\longmapsto s_{u(\alpha)}
\]
extends to an isomorphism
\[
 u:W(\cR)\xrightarrow{\sim}W(\cR')
\]
such that
\[
 u(w(\alpha))=u(w)\bigl(u(\alpha)\bigr).
\]
\end{lemma}

\begin{proof}
It suffices to prove, for \(\alpha,\beta\in R\), that
\[
 u(s_\alpha(\beta))=s_{u(\alpha)}u(\beta),
 \qquad
 q(s_\alpha(\beta))=q(\beta).
\]
Successively, one has
\begin{align*}
 f\bigl(s_{u(\alpha)}u(\beta)\bigr)
  &=f(u(\beta))
    -\bigl(u(\alpha)^\ast,u(\beta)\bigr)f(u(\alpha))\\
  &=q(\beta)\beta
    -q(\beta)q(\alpha)^{-1}
       (\alpha^\ast,\beta)q(\alpha)\alpha\\
  &=q(\beta)\bigl(\beta-(\alpha^\ast,\beta)\alpha\bigr)\\
  &=q(\beta)s_\alpha(\beta).
\end{align*}
If
\[
 \gamma=u^{-1}\bigl(s_{u(\alpha)}u(\beta)\bigr),
\]
then
\[
 q(\gamma)\gamma=f(u(\gamma))=q(\beta)s_\alpha(\beta).
\]
The roots \(\gamma\) and \(s_\alpha(\beta)\) are proportional over
\(\mathbb Q\), hence equal or opposite. Since \(q(\gamma)\) and
\(q(\beta)\) are positive, it follows that
\[
 q(\gamma)=q(\beta),
 \qquad
 \gamma=s_\alpha(\beta).
\]
\end{proof}

\begin{definition}\label{sga3:XXI:6.8.5}
The numbers \(q(\alpha)\) are called the \emph{root exponents} of
\(f\).
\end{definition}

\begin{example}\label{sga3:XXI:6.8.6}
Let \(\cR\) be a reduced root datum and let
\(q=p^n\), \(n\in\mathbb N\). Multiplication by \(q\),
\[
 M\longrightarrow M,\qquad x\longmapsto qx,
\]
is a \(p\)-morphism whose root exponents are all equal to \(q\)
(and for which \(u=\operatorname{id}\)). It is denoted
\[
 q:\cR\longrightarrow\cR.
\]
\end{example}
